\section{Existence by generators and relations}
\label{sec:fixed-base-existence}

Drieux's 2026 slides propose computing indexed colimits from cocompleteness,
change-of-base adjoints, and Beck--Chevalley
\cite[logical slides~17--18]{Drieux2026TACL}.  The slides do not give a
published proof of the localized coend defined in
\cref{def:relative-cowedge}.  We therefore begin with a fixed base object and
derive the representing object directly.  The result is deliberately stated
with only the colimits actually used.

\subsection{A mate calculation}

Assume for the moment that every reindexing functor
\(v^*:\V(U)\to\V(V)\) has a chosen left adjoint
\(v_!\dashv v^*\), with unit \(\eta^v\) and counit \(\varepsilon^v\).
For composable arrows
\[
  W\xrightarrow{w}V\xrightarrow{v}U,
\]
both \((vw)_!\) and \(v_!w_!\) are left adjoint, up to the constraint
\(\alpha_{w,v}:w^*v^*\cong(vw)^*\), to the same functor.  Uniqueness of
left adjoints gives a canonical natural isomorphism
\[
  c_{v,w}:(vw)_!\xrightarrow{\ \cong\ }v_!w_!.
\]
Its direction and normalization are fixed by requiring precomposition with
\(c_{v,w}\) to be the natural chain
\[
\begin{split}
  \V(U)(v_!w_!Z,X)
  &\cong\V(W)(Z,w^*v^*X)\\
  &\cong\V(W)(Z,(vw)^*X)
  \cong\V(U)((vw)_!Z,X).
\end{split}
\]
The uniqueness part of the adjunction correspondence makes the \(c_{v,w}\)
coherent for triples; this is the usual composition constraint obtained by
transporting composite adjunctions.
For \(Z\in\V(V)\), put
\begin{equation}
  \beta_{v,w,Z}:
  (vw)_!w^*Z\xrightarrow{c_{v,w}}
  v_!w_!w^*Z\xrightarrow{v_!\varepsilon^w_Z}v_!Z.
  \label{eq:beta-mate}
\end{equation}

\begin{lemma}[The base-coherence mate]
\label{lem:base-mate}
Let \(h:Z\to v^*X\), and let
\(\bar h:v_!Z\to X\) be its adjunct.  Under
\((vw)_!\dashv(vw)^*\), the adjunct of
\[
  w^*Z\xrightarrow{w^*h}w^*v^*X
  \xrightarrow{\alpha_{w,v,X}}(vw)^*X
\]
is the composite
\[
  (vw)_!w^*Z\xrightarrow{\beta_{v,w,Z}}v_!Z
  \xrightarrow{\bar h}X.
\]
\end{lemma}

\begin{proof}
The adjunct of \(h\) is
\(\bar h=\varepsilon^v_X\,v_!h\).  Hence the second displayed composite is
\[
  (vw)_!w^*Z
  \xrightarrow{c_{v,w}}v_!w_!w^*Z
  \xrightarrow{v_!\varepsilon^w_Z}v_!Z
  \xrightarrow{v_!h}v_!v^*X
  \xrightarrow{\varepsilon^v_X}X.
\]
Naturality of \(\varepsilon^w\) moves \(w_!w^*h\) across
\(\varepsilon^w_Z\).  The result is the counit formula for the composite
adjunction \(v_!w_!\dashv w^*v^*\).  Transporting that composite adjunction
along \(c_{v,w}\) and \(\alpha_{w,v}\) gives the counit formula for
\((vw)_!\dashv(vw)^*\).  Taking its adjunct therefore yields
\(\alpha_{w,v,X}w^*h\), as asserted.  This is precisely the standard
calculus of mates; the displayed computation fixes the direction of every
constraint used later.
\end{proof}

\subsection{The coequalizer presentation}

Fix \(U\in\C\).  Since \(\C\) and every \(\D(V)\) are small, the following
indexing collections are sets.  Define the object of generators
\begin{equation}
  P_U=\coprod_{\substack{v:V\to U\\A\in\Ob\D(V)}}
      v_!H_V(A,A),
  \label{eq:PU}
\end{equation}
and write
\(\iota_{v,A}:v_!H_V(A,A)\to P_U\) for its coproduct injections.

There are two sorts of relations.  Fibrewise dinaturality is generated by
\begin{equation}
  R_U^{\mathrm{din}}=
  \coprod_{\substack{v:V\to U\\f:A\to B\text{ in }\D(V)}}
      v_!H_V(B,A).
  \label{eq:Rdin}
\end{equation}
On the summand indexed by \((v,f:A\to B)\), define the parallel maps to
\(P_U\) by
\begin{align}
  d^{\mathrm{din}}_0
  &=\iota_{v,A}\,v_!H_V(f^{\op},1_A),
  \\
  d^{\mathrm{din}}_1
  &=\iota_{v,B}\,v_!H_V(1_B,f).
  \label{eq:din-relation-maps}
\end{align}

Base coherence is generated by
\begin{equation}
  R_U^{\mathrm{bas}}=
  \coprod_{\substack{v:V\to U,\ w:W\to V\\A\in\Ob\D(V)}}
      (vw)_!w^*H_V(A,A).
  \label{eq:Rbase}
\end{equation}
On the summand indexed by \((v,w,A)\), its two maps to \(P_U\) are
\begin{align}
  d^{\mathrm{bas}}_0
  &=(vw)_!w^*H_V(A,A)
    \xrightarrow{(vw)_!\chi^H_{w,A,A}}
    (vw)_!H_W(w^*A,w^*A)
    \xrightarrow{\iota_{vw,w^*A}}P_U,
  \label{eq:base-relation-first}
  \\
  d^{\mathrm{bas}}_1
  &=(vw)_!w^*H_V(A,A)
    \xrightarrow{\beta_{v,w,H_V(A,A)}}v_!H_V(A,A)
    \xrightarrow{\iota_{v,A}}P_U.
  \label{eq:base-relation-second}
\end{align}
Let
\begin{equation}
  R_U=R_U^{\mathrm{din}}\amalg R_U^{\mathrm{bas}},
  \qquad
  d_i=d_i^{\mathrm{din}}\amalg d_i^{\mathrm{bas}}\quad(i=0,1).
  \label{eq:Rtotal}
\end{equation}

\begin{theorem}[Fixed-base coequalizer formula]
\label{thm:fixed-base-coequalizer}
Suppose that every \(v^*\) has a chosen left adjoint \(v_!\), and that
\(\V(U)\) has the coproducts in
\cref{eq:PU,eq:Rdin,eq:Rbase}, the binary coproduct in
\cref{eq:Rtotal}, and the coequalizer
\begin{equation}
  R_U\mathrel{\substack{\xrightarrow{d_0}\\[-0.35em]
                         \xrightarrow[d_1]{}}}
  P_U\xrightarrow{q_U}L_U.
  \label{eq:relative-coend-coeq}
\end{equation}
Then \(L_U\), with the cowedge defined below, is the relative coend of
\(H\) over \(U\).  In particular,
\[
  \V(U)(L_U,X)\cong\RCow_U(H;X)
\]
naturally in \(X\in\V(U)\).

No pullbacks in \(\C\), Beck--Chevalley isomorphisms, or preservation of
colimits by reindexing are needed for this fixed-base statement.
\end{theorem}

\begin{proof}
We give the construction and both directions of the representing bijection.

\emph{Step 1: the universal components.}
For each \(v:V\to U\) and \(A\in\D(V)\), transpose
\[
  v_!H_V(A,A)\xrightarrow{\iota_{v,A}}P_U
  \xrightarrow{q_U}L_U
\]
through \(v_!\dashv v^*\).  Denote the resulting map by
\begin{equation}
  \lambda^U_{v,A}:H_V(A,A)\longrightarrow v^*L_U.
  \label{eq:coeq-universal-component}
\end{equation}

\emph{Step 2: fibrewise dinaturality.}
For \(f:A\to B\) in \(\D(V)\), the equation
\(q_Ud^{\mathrm{din}}_0=q_Ud^{\mathrm{din}}_1\) on the
\((v,f)\)-summand says
\[
  q_U\iota_{v,A}v_!H_V(f^{\op},1_A)
  =q_U\iota_{v,B}v_!H_V(1_B,f).
\]
Taking adjuncts under \(v_!\dashv v^*\) gives exactly
\[
  \lambda^U_{v,A}H_V(f^{\op},1_A)
  =\lambda^U_{v,B}H_V(1_B,f).
\]

\emph{Step 3: base coherence.}
On the \((v,w,A)\)-summand of \(R_U^{\mathrm{bas}}\), the equation
\(q_Ud^{\mathrm{bas}}_0=q_Ud^{\mathrm{bas}}_1\) has two adjuncts under
\((vw)_!\dashv(vw)^*\).  The adjunct of the first is
\[
  w^*H_V(A,A)\xrightarrow{\chi^H_{w,A,A}}
  H_W(w^*A,w^*A)
  \xrightarrow{\lambda^U_{vw,w^*A}}(vw)^*L_U.
\]
By \cref{lem:base-mate}, the adjunct of the second is
\[
  w^*H_V(A,A)\xrightarrow{w^*\lambda^U_{v,A}}
  w^*v^*L_U\xrightarrow{\alpha_{w,v,L_U}}(vw)^*L_U.
\]
Their equality is precisely \cref{eq:base-coherence}.  Thus
\(\lambda^U\) is a relative cowedge.

\emph{Step 4: maps out of the coequalizer give cowedges.}
Given \(m:L_U\to X\), define
\[
  \omega^m_{v,A}=v^*m\,\lambda^U_{v,A}.
\]
The equations already proved for \(\lambda^U\), together with naturality of
the pseudofunctor constraints, show that \(\omega^m\) is a relative cowedge.

\emph{Step 5: every cowedge factors uniquely.}
Conversely, let \(\omega\in\RCow_U(H;X)\).  Transpose each component to
\[
  \bar\omega_{v,A}:v_!H_V(A,A)\longrightarrow X.
\]
The coproduct property gives a unique
\(\bar\omega:P_U\to X\) with
\(\bar\omega\iota_{v,A}=\bar\omega_{v,A}\).  On a summand of
\(R_U^{\mathrm{din}}\), the equality
\(\bar\omega d_0=\bar\omega d_1\) is, after adjunction, exactly the
fibrewise dinaturality of \(\omega\).  On a summand of
\(R_U^{\mathrm{bas}}\), \cref{lem:base-mate} identifies the same equality
with the base-coherence equation for \(\omega\).  Thus \(\bar\omega\)
coequalizes \(d_0,d_1\), so there is a unique \(m:L_U\to X\) such that
\(mq_U=\bar\omega\).  Taking adjuncts on each generator shows
\(v^*m\,\lambda^U_{v,A}=\omega_{v,A}\).

The two assignments are inverse because both are determined on the coproduct
generators, and they are natural in \(X\) because postcomposition is used at
every step.  \Cref{prop:representing} now identifies \(L_U\) as the relative
coend.
\end{proof}

\begin{remark}[Relation to the ordinary formula]
\label{rem:ordinary-coeq-boundary}
When \(\C=\one\), the base-relation summands only encode unit constraints,
and \cref{eq:relative-coend-coeq} reduces to the standard coequalizer by
objects and arrows of \(\D(*)\).  This agrees with the ordinary definition as
an initial cowedge \cite[Definitions~1.1.4 and 1.1.6,
pp.~23--24]{Loregian2023}.  On a nonterminal base, the summand
\(R_U^{\mathrm{bas}}\) is additional data; deleting it changes the universal
problem.
\end{remark}


\section{Base-change stability and the missing hypothesis}
\label{sec:stability-criterion}

The fixed-base construction does not by itself imply that the maps
\(b_u:L_{U'}\to u^*L_U\) are invertible.  The exact criterion is immediate
from the universal property.

\begin{proposition}[Exact stability criterion]
\label{prop:stability-criterion}
For \(u:U'\to U\), assume that relative coends exist over \(U\) and
\(U'\).  The canonical map \(b_u\) is an isomorphism if and only if the
restricted universal cowedge
\[
  \rho_u(L_U,\lambda^U)=(u^*L_U,\rho_u\lambda^U)
\]
is initial in \(\Cow_{U'}(H)\).
\end{proposition}

\begin{proof}
If \(b_u\) is invertible, it is an isomorphism in \(\Cow_{U'}(H)\) from the
initial cowedge \((L_{U'},\lambda^{U'})\) to the restricted cowedge; hence
the latter is initial.  Conversely, if both cowedges are initial, there is a
unique arrow in each direction between them.  The two composites are
endomorphisms of initial objects and therefore identities.  The arrow from
\((L_{U'},\lambda^{U'})\) is, by its defining equation, \(b_u\); hence it is
invertible.
\end{proof}

The following example isolates a genuine gap in a tempting statement of the
hypotheses.

\begin{example}[Available Beck--Chevalley squares are not enough]
\label{ex:missing-pullback}
Let \(\C\) have three objects and two nonidentity arrows
\[
  0\xrightarrow{p}2\xleftarrow{q}1,
\]
and equip it with the trivial Grothendieck topology.  Let \(\V\) be the
constant indexed category \(\Set\), so every reindexing functor and every
chosen left adjoint is the identity.  All coproducts and coequalizers exist
and are preserved, and the Beck--Chevalley map for every pullback square that
exists is an identity.

Put
\[
  \D(2)=\varnothing,
  \qquad
  \D(0)=\D(1)=\one,
\]
with the unique reindexing functors, and set
\(H_0(*,*)=H_1(*,*)=\one\).  There is no component over \(2\).

A relative cowedge over \(2\) with apex \(X\) consists of one point of
\(X\) contributed by \((p,*)\) and one independently contributed by
\((q,*)\).  There is no arrow in a slice relating these two components.
Consequently
\[
  \RCow_2(H;X)\cong X\times X
  \quad\text{and}\quad
  L_2\cong\{0,1\}.
\]
Over \(0\) and \(1\) there is one component, so
\(L_0\cong L_1\cong\one\).  Since reindexing in \(\V\) is the identity,
\[
  b_p:\one\longrightarrow\{0,1\},
  \qquad
  b_q:\one\longrightarrow\{0,1\}
\]
are the two coproduct injections and are not isomorphisms.

Thus cocomplete fibres, cocontinuous reindexing, left adjoints, and
Beck--Chevalley for all \emph{available} pullback squares do not imply
stability on an arbitrary site.  What is missing here is the pullback of
\(p\) and \(q\).  The next theorem shows that chosen pullbacks, together with
the expected exactness, close precisely this gap for the displayed
presentation.
\end{example}

\begin{lemma}[Pasting for the canonical Beck--Chevalley mates]
\label{lem:BC-pasting}
For a chosen pullback square as in \cref{eq:BC-square}, let
\[
  \mathrm{BC}_{u,v}:s_!r^*\Longrightarrow u^*v_!
\]
be the left mate of the pseudofunctorial comparison
\(r^*v^*\cong s^*u^*\).  Given \(w:T\to V\), form the pullback of
\(vw\) along \(u\),
\[
  Z\xrightarrow{k}T,\qquad Z\xrightarrow{t}U',
\]
and let \(\bar w:Z\to W\) be determined by
\(r\bar w=wk\) and \(s\bar w=t\).  Then, with the canonical
pseudofunctor constraints suppressed, use the following convention: if the
induced square \((k,\bar w;r,w)\) is not the selected pullback of \(w\)
along \(r\), transport its canonical mate along the unique isomorphism to
the selected pullback.  Canonical mates and their invertibility are
unchanged by this transport.  With that convention,
\begin{equation}
\begin{split}
  \mathrm{BC}_{u,vw}
  ={}&u^*(c_{v,w}^{-1})\,
      (\mathrm{BC}_{u,v}w_!)\,
      s_!(\mathrm{BC}_{r,w})\,
      c_{s,\bar w}.
\end{split}
\label{eq:BC-pasting}
\end{equation}
Moreover, for \(Z_0\in\V(V)\),
\begin{equation}
\begin{split}
 u^*\beta_{v,w,Z_0}\,
 \mathrm{BC}_{u,vw,w^*Z_0}
 ={}&
 \mathrm{BC}_{u,v,Z_0}\,
 \beta_{s,\bar w,r^*Z_0}\,
 t_!(\zeta_{Z_0}),
\end{split}
\label{eq:BC-beta}
\end{equation}
where
\(\zeta_{Z_0}:k^*w^*Z_0\cong\bar w^*r^*Z_0\) is induced by
\(wk=r\bar w\).
\end{lemma}

\begin{proof}
Take right mates of both sides of \cref{eq:BC-pasting}, first under
\(t_!\dashv t^*\) and then under \((vw)_!\dashv(vw)^*\).  The right mate of
the left-hand side is, by definition, the comparison associated to the
outer rectangle
\[
  k^*w^*v^*\cong t^*u^*.
\]
On the right-hand side, the mates of
\(\mathrm{BC}_{r,w}\) and \(\mathrm{BC}_{u,v}\) are the comparisons for
the two constituent pullback squares, while the mates of
\(c_{s,\bar w}\) and \(c^{-1}_{v,w}\) merely reassociate the two composite
adjunctions.  After the two triangle identities cancel the inserted units
and counits, the result is the pasted comparison
\[
  k^*w^*v^*
  \cong\bar w^*r^*v^*
  \cong\bar w^*s^*u^*
  \cong t^*u^*.
\]
The pseudofunctor pentagon says that this is the outer-rectangle comparison.
Since taking mates is bijective, \cref{eq:BC-pasting} follows.

For \cref{eq:BC-beta}, expand each \(\beta\) by
\cref{eq:beta-mate}, substitute \cref{eq:BC-pasting}, and use naturality of
\(\mathrm{BC}_{u,v}\) with respect to the counit
\(\varepsilon^w:w_!w^*\to1\).  The middle composite is then
\[
  s_!\bar w_!k^*w^*Z_0
  \xrightarrow{s_!\bar w_!(\zeta_{Z_0})}
  s_!\bar w_!\bar w^*r^*Z_0
  \xrightarrow{s_!\varepsilon^{\bar w}}s_!r^*Z_0,
\]
on both sides.  The remaining \(c\)'s agree by the triple coherence proved
from uniqueness of composite adjoints above.  This gives
\cref{eq:BC-beta}.  The analogous replacements involving \(\chi^H\) agree
by the composition axiom for the indexed functor \(H\); that is the
compatibility used on base-relation summands below.
\end{proof}

\begin{theorem}[Stable base change from pullbacks]
\label{thm:stable-pullback}
Assume the hypotheses of \cref{thm:fixed-base-coequalizer} for every base
object.  Assume in addition that:
\begin{enumerate}[label=(S\arabic*)]
  \item \(\C\) has chosen pullbacks;
  \item the canonical mates of the chosen adjoints satisfy
  Beck--Chevalley: for every chosen
  pullback square
  \begin{equation}
  \begin{tikzcd}[column sep=3.6em,row sep=2.4em]
    W \arrow[r,"r"] \arrow[d,"s"'] & V \arrow[d,"v"]\\
    U' \arrow[r,"u"'] & U
  \end{tikzcd}
  \qquad
  \mathrm{BC}_{u,v}:s_!r^*\xrightarrow{\ \cong\ }u^*v_!;
  \label{eq:BC-square}
  \end{equation}
  \item every \(u^*\) preserves the coproducts and coequalizers appearing in
        \cref{eq:PU,eq:Rdin,eq:Rbase,eq:Rtotal,eq:relative-coend-coeq},
        including the binary coproduct
        \(R_U^{\mathrm{din}}\amalg R_U^{\mathrm{bas}}\).
\end{enumerate}
Then every canonical base-change map
\[
  b_u:L_{U'}\xrightarrow{\ \cong\ }u^*L_U
\]
is an isomorphism.  Hence the fixed-base coequalizers assemble, uniquely, to
a cartesian pseudonatural object of \(\V\).
\end{theorem}

\begin{proof}
Fix \(u:U'\to U\).  We construct an inverse
\(a_u:u^*L_U\to L_{U'}\) on generators and then check both composites.
Write
\[
  j^U_{v,A}=q_U\iota_{v,A}:v_!H_V(A,A)\longrightarrow L_U.
\]

\emph{Step 1: pull a generator across the square.}
For every \(v:V\to U\), choose its pullback along \(u\) as in
\cref{eq:BC-square}.  Preservation of the coproduct in \cref{eq:PU},
Beck--Chevalley, and the indexed comparison of \(H\) give
\begin{align}
  u^*P_U
  &\cong
  \coprod_{v:V\to U,\ A\in\D(V)}u^*v_!H_V(A,A)
  \\
  &\xrightarrow{\coprod\mathrm{BC}_{u,v}^{-1}}
  \coprod_{v,A}s_!r^*H_V(A,A)
  \\
  &\xrightarrow{\coprod s_!\chi^H_{r,A,A}}
  \coprod_{v,A}s_!H_W(r^*A,r^*A)
  \xrightarrow{\coprod\iota_{s,r^*A}}P_{U'}
  \xrightarrow{q_{U'}}L_{U'}.
  \label{eq:a-on-generators}
\end{align}
Call this composite \(a_P:u^*P_U\to L_{U'}\).
Consequently, once \(a_u\) is induced from \(a_P\), its defining generator
equation is
\begin{equation}
  a_u\,u^*j^U_{v,A}
  =
  j^{U'}_{s,r^*A}\,
  s_!\chi^H_{r,A,A}\,
  \mathrm{BC}_{u,v,H_V(A,A)}^{-1}.
  \label{eq:au-generator}
\end{equation}

\emph{Step 2: the relations are respected.}
For a dinaturality summand indexed by \((v,f:A\to B)\), pullback turns
\(f\) into \(r^*f:r^*A\to r^*B\).  Under the first three isomorphisms in
\cref{eq:a-on-generators}, the two maps are exactly the two legs of the
dinaturality relation indexed by \((s,r^*f)\) in
\(R_{U'}^{\mathrm{din}}\).  They therefore agree after \(q_{U'}\).

For a base-relation summand indexed by
\(v:V\to U\), \(w:T\to V\), and \(A\in\D(V)\), form the chosen pullbacks
\[
  W=U'\times_UV,
  \qquad
  Z=U'\times_UT.
\]
Writing \((s,r):W\to U'\times V\) and
\((t,k):Z\to U'\times T\), the pullback universal property gives a unique
\(\bar w:Z\to W\) with
\[
  s\bar w=t,
  \qquad
  r\bar w=wk.
\]
Pasting of Beck--Chevalley squares identifies the pullback of the first leg
\cref{eq:base-relation-first} with the generator
\((t,\bar w^*r^*A)\), and identifies the pullback of the second leg
\cref{eq:base-relation-second} with the generator \((s,r^*A)\).  The
comparison
\(\bar w^*r^*A\cong k^*w^*A\) is the pseudofunctorial constraint of
\(\D\); the analogous comparison for \(H\) is the coherence of \(\chi^H\).
Precisely, \cref{eq:BC-pasting} gives the comparison for the first leg,
while \cref{eq:BC-beta} gives the comparison for the \(\beta\)-map in the
second leg.  The composition axiom for \(\chi^H\) transports both to the
same chosen identification of
\(\bar w^*r^*H_V(A,A)\) with \(k^*w^*H_V(A,A)\).
More explicitly, after these identifications the two composites into
\(L_{U'}\) are
\begin{align*}
  t_!\bar w^*H_W(r^*A,r^*A)
  &\xrightarrow{t_!\chi^H_{\bar w,r^*A,r^*A}}
  t_!H_Z(\bar w^*r^*A,\bar w^*r^*A)
  \xrightarrow{q_{U'}\iota_{t,\bar w^*r^*A}}L_{U'},\\
  t_!\bar w^*H_W(r^*A,r^*A)
  &\xrightarrow{\beta_{s,\bar w,H_W(r^*A,r^*A)}}
  s_!H_W(r^*A,r^*A)
  \xrightarrow{q_{U'}\iota_{s,r^*A}}L_{U'}.
\end{align*}
They are precisely the two legs of the base relation in
\(R_{U'}^{\mathrm{bas}}\) indexed by \((s,\bar w,r^*A)\).  The second leg
has the displayed direction by \cref{lem:base-mate}.  Thus it too is
equalized by \(q_{U'}\).

It follows that \(a_P\) coequalizes \(u^*d_0,u^*d_1\).  Since \(u^*\)
preserves \cref{eq:relative-coend-coeq}, there is a unique
\begin{equation}
  a_u:u^*L_U\longrightarrow L_{U'}
  \quad\text{with}\quad
  a_u\,u^*q_U=a_P.
  \label{eq:au}
\end{equation}

We shall also use the following generator form of the defining equation for
\(b_u\):
\begin{equation}
  b_u\,j^{U'}_{s,r^*A}\,s_!\chi^H_{r,A,A}
  =
  u^*j^U_{v,A}\,\mathrm{BC}_{u,v,H_V(A,A)}.
  \label{eq:bu-generator}
\end{equation}
Both sides have domain \(s_!r^*H_V(A,A)\).  To verify the equality, take
adjuncts under \(s_!\dashv s^*\).  The left adjunct, using the defining
identity \(b_u\lambda^{U'}=\rho_u\lambda^U\), is the composite represented
by
\[
  r^*H_V(A,A)\xrightarrow{\chi^H_{r,A,A}}
  H_W(r^*A,r^*A)\xrightarrow{\lambda^U_{vr,r^*A}}
  (vr)^*L_U\cong s^*u^*L_U.
\]
The right adjunct is the same composite: this is the definition of the
canonical mate \(\mathrm{BC}_{u,v}\), followed by the base-coherence
equation for \(\lambda^U\) along \(r\).  Adjunction is bijective, proving
\cref{eq:bu-generator}.

\emph{Step 3: \(a_ub_u=1\).}
It suffices to precompose with every coproduct generator
\(v'_!H_{V'}(A',A')\to P_{U'}\to L_{U'}\), since a coequalizer is an
epimorphism.  Pull back the composite \(uv':V'\to U\) along \(u\):
\[
  W=U'\times_UV'
  \xrightarrow{r}V',
  \qquad
  W\xrightarrow{s}U'.
\]
There is a canonical section
\(\delta=(v',1_{V'}):V'\to W\) with
\(s\delta=v'\) and \(r\delta=1_{V'}\).  By the definition of \(b_u\), the
following section mate makes this statement explicit.  Put
\(Z_0=H_{V'}(A',A')\), and let
\[
  \widehat\alpha_{\delta,r,Z_0}:\delta_!Z_0\longrightarrow r^*Z_0
\]
be the adjunct of the inverse of the constraint
\(\delta^*r^*Z_0\cong(r\delta)^*Z_0\cong Z_0\).  Define
\begin{equation}
  \Gamma_{\delta,A'}:
  v'_!Z_0\cong(s\delta)_!Z_0
  \xrightarrow{c_{s,\delta}}s_!\delta_!Z_0
  \xrightarrow{s_!\widehat\alpha_{\delta,r,Z_0}}s_!r^*Z_0.
  \label{eq:section-mate}
\end{equation}
The base relation of \(L_{U'}\) indexed by
\((s,\delta,r^*A')\), together with \(r\delta=1\) and the composition
coherence of \(\chi^H\), is exactly
\begin{equation}
  j^{U'}_{s,r^*A'}\,s_!\chi^H_{r,A',A'}\,\Gamma_{\delta,A'}
  =j^{U'}_{v',A'}.
  \label{eq:section-base-relation}
\end{equation}
Precomposing \cref{eq:bu-generator}, for the pullback of \(uv'\) along
\(u\), with \(\Gamma_{\delta,A'}\), and then using
\cref{eq:section-base-relation}, gives
\begin{equation}
  b_u j^{U'}_{v',A'}
  =
  u^*j^U_{uv',A'}\,
  \mathrm{BC}_{u,uv',Z_0}\,
  \Gamma_{\delta,A'}.
  \label{eq:bu-section}
\end{equation}
Now \cref{eq:au-generator,eq:section-base-relation,eq:bu-section} give the
calculation
\[
\begin{split}
  a_ub_u j^{U'}_{v',A'}
  &=
  a_u u^*j^U_{uv',A'}\,
  \mathrm{BC}_{u,uv',Z_0}\,\Gamma_{\delta,A'}\\
  &=
  j^{U'}_{s,r^*A'}s_!\chi^H_{r,A',A'}\,
  \Gamma_{\delta,A'}\\
  &=j^{U'}_{v',A'}.
\end{split}
\]
The maps \(j^{U'}_{v',A'}\) are jointly epimorphic because they are the
coproduct injections followed by the coequalizer.  Hence \(a_ub_u=1\).

\emph{Step 4: \(b_ua_u=1\).}
Because \(u^*\) preserves the coequalizer, \(u^*q_U\) is an epimorphism.
Together with preservation of \cref{eq:PU}, it is therefore enough to
precompose with every \(u^*j^U_{v,A}\).  Equations
\cref{eq:au-generator,eq:bu-generator} give, on the pullback square for
\(u\) and \(v\),
\[
\begin{split}
  b_ua_u\,u^*j^U_{v,A}
  &=
  b_u j^{U'}_{s,r^*A}s_!\chi^H_{r,A,A}
       \mathrm{BC}_{u,v,H_V(A,A)}^{-1}\\
  &=
  u^*j^U_{v,A}\,
       \mathrm{BC}_{u,v,H_V(A,A)}
       \mathrm{BC}_{u,v,H_V(A,A)}^{-1}\\
  &=u^*j^U_{v,A}.
\end{split}
\]
Thus \(b_ua_u=1_{u^*L_U}\).

Thus \(a_u=b_u^{-1}\).  Coherent assembly now follows from
\cref{thm:assembly}.
\end{proof}

\begin{remark}[Exact comparison with Shulman's hypotheses]
\label{rem:shulman-boundary}
Shulman calls the existence of \(f_!\dashv f^*\) together with
Beck--Chevalley for pullback squares ``indexed coproducts''
\cite[Definition~2.4, p.~7]{Shulman2013}.  His tensor-preservation condition
is Definition~2.10, and indexed profunctor composition is the coequalizer of
Lemma~3.25 \cite[pp.~8 and 21--22]{Shulman2013}.  Theorem
\ref{thm:stable-pullback} is not a restatement of that lemma: our
\(R_U^{\mathrm{bas}}\) explicitly imposes relations between all generalized
objects in \(\C/U\), while Shulman's formula composes enriched indexed
profunctors.  A comparison between the resulting objects still requires a
typed identification of these two universal problems.
\end{remark}


\section{A Fubini theorem for stable relative coends}
\label{sec:relative-fubini}

The ordinary Fubini theorem identifies a coend over a product category with
either order of iterated coends
\cite[Theorem~1.3.1, pp.~37--39]{Loregian2023}.  In the relative setting an
inner coend must itself vary as indexed data before an outer relative coend
can even be typed.  This is why stability, rather than mere fibrewise
existence, is the correct hypothesis.

The indexed use of Fubini is not itself new: Pitts invokes it in the
\(S\)-indexed setting in the proof of his product theorem
\cite[proof of Theorem~2.3, p.~51]{Pitts1985}.  More recently, Nickel
reported joint work with Arkor on coends in compact closed virtual
equipments, explicitly aiming to recover Fubini and to include
\(\mathbb{S}\mathsf{pan}(\mathcal E)\) as an example
\cite{NickelArkor2026}.  The available source for the latter is a seminar
abstract, not a public manuscript.  The result below addresses a different
fixed-site question: when a full-slice inner universal object is sufficiently
stable to serve as the input of another full-slice coend.

Let
\[
  \mathbb A,\mathbb B:\C^{\op}\longrightarrow\Cat
\]
be small-fibred indexed categories, and let
\begin{equation}
  H:\mathbb A^{\vee}\times_{\C}\mathbb A
       \times_{\C}\mathbb B^{\vee}\times_{\C}\mathbb B
       \longrightarrow\V
  \label{eq:four-variable-H}
\end{equation}
be an indexed functor.  Put
\[
  \widehat H_V\bigl((A',B'),(A,B)\bigr)
  =H_V(A',A,B',B).
\]
Then \(\widehat H\) is an indexed functor on
\((\mathbb A\times_{\C}\mathbb B)^{\vee}
 \times_{\C}(\mathbb A\times_{\C}\mathbb B)\).

\begin{definition}[Simultaneous relative cowedge]
\label{def:simultaneous-cowedge}
For \(U\in\C\) and \(X\in\V(U)\), a simultaneous relative cowedge for
\(H\) is a family
\[
  \omega_{v,A,B}:H_V(A,A,B,B)\longrightarrow v^*X
\]
for \(v:V\to U\), \(A\in\mathbb A(V)\), and
\(B\in\mathbb B(V)\), which is dinatural separately in \(A\) and in
\(B\).  Explicitly, for \(f:A_0\to A_1\) and \(g:B_0\to B_1\),
\begin{align}
 \omega_{v,A_0,B}\,
 H_V(f^{\op},1_{A_0},1_B,1_B)
 &=
 \omega_{v,A_1,B}\,
 H_V(1_{A_1},f,1_B,1_B),
 \label{eq:simultaneous-A-dinaturality}\\
 \omega_{v,A,B_0}\,
 H_V(1_A,1_A,g^{\op},1_{B_0})
 &=
 \omega_{v,A,B_1}\,
 H_V(1_A,1_A,1_{B_1},g).
 \label{eq:simultaneous-B-dinaturality}
\end{align}
The common domains are respectively
\(H_V(A_1,A_0,B,B)\) and \(H_V(A,A,B_1,B_0)\).
Finally, the family satisfies the joint base-coherence equation
\[
  \omega_{vw,w^*A,w^*B}\,\chi^H_{w,A,A,B,B}
  =\alpha_{w,v,X}\,w^*\omega_{v,A,B}.
\]
Its set is denoted \(\RBiCow_U(H;X)\).
\end{definition}

\begin{lemma}[Product versus separate dinaturality]
\label{lem:product-dinaturality}
There is a natural equality of cowedge sets
\[
  \RCow_U(\widehat H;X)
  =\RBiCow_U(H;X).
\]
\end{lemma}

\begin{proof}
Dinaturality for a morphism \((f,g):(A,B)\to(A',B')\) in a product fibre
implies the two separate equations by taking \(g\) or \(f\) to be an
identity.  Conversely, factor
\[
  (f,g)=(1_{A'},g)(f,1_B).
\]
Apply \(A\)-dinaturality to the first change of object and
\(B\)-dinaturality to the second; functoriality of \(H_V\) makes the middle
terms agree.  This gives dinaturality for \((f,g)\).  The base-coherence
equation is literally the same for \(\widehat H\), so the two sets coincide,
naturally in \(X\).
\end{proof}

For \(B',B\in\mathbb B(U)\), restrict \(H\) to the slice \(\C/U\) by
setting
\[
  H^{B',B}_w(A',A)
  =H_W(A',A,w^*B',w^*B)
  \qquad(w:W\to U).
\]
Assume that its relative coend exists and write
\begin{equation}
  K_U(B',B)=
  \mathop{\int^{A\in\mathbb A}_{\Rel,U}}
  H(A,A,B',B).
  \label{eq:inner-K}
\end{equation}
Its universal components are
\begin{equation}
  \kappa^{U,B',B}_{w,A}:
  H_W(A,A,w^*B',w^*B)\longrightarrow w^*K_U(B',B).
  \label{eq:kappa-components}
\end{equation}
Whenever the inner coends over \(U\) and \(U'\) exist, write
\begin{equation}
  b_u^{B',B}:
  K_{U'}(u^*B',u^*B)\longrightarrow u^*K_U(B',B)
  \label{eq:K-base-map}
\end{equation}
for their canonical base-change map from \cref{eq:base-change-map}.

\begin{lemma}[The inner universal cowedge under base change]
\label{lem:kappa-base-change}
For \(u:U'\to U\), \(w:W\to U'\), \(A\in\mathbb A(W)\), and
\(B',B\in\mathbb B(U)\), the canonical map
\cref{eq:K-base-map} is characterized by
\begin{equation}
\begin{split}
 w^*b_u^{B',B}\,
 \kappa^{U',u^*B',u^*B}_{w,A}
 ={}&
 \alpha^{-1}_{w,u,K_U(B',B)}\,
 \kappa^{U,B',B}_{uw,A}\\
 &{}\quad
 H_W\bigl(1,1,
   (\alpha^{-1}_{w,u,B'})^{\op},
   \alpha_{w,u,B}\bigr).
\end{split}
\label{eq:kappa-base-change}
\end{equation}
The suppressed unit constraints are the same on both sides.
\end{lemma}

\begin{proof}
The right-hand side is the component at \((w,A)\) of the restriction along
\(u\) of the universal inner cowedge over \(U\): the last arrow transports
the parameters \(w^*u^*B',w^*u^*B\) to
\((uw)^*B',(uw)^*B\), and the preceding \(\alpha^{-1}\) transports its
apex to \(w^*u^*K_U(B',B)\).  By definition,
\(b_u^{B',B}\) is the unique map from the universal cowedge over \(U'\) to
this restricted cowedge.  Its component equation is therefore exactly
\cref{eq:kappa-base-change}.
\end{proof}

\begin{lemma}[Stable inner coends form an indexed bifunctor]
\label{lem:inner-indexed-functor}
Suppose the objects in \cref{eq:inner-K} exist for every \(U,B',B\), and
the canonical base-change maps in \cref{eq:K-base-map} are isomorphisms.
Then the \(K_U\) assemble to an indexed functor
\[
  K:\mathbb B^{\vee}\times_{\C}\mathbb B\longrightarrow\V,
\]
whose comparison along \(u\) is \((b_u^{B',B})^{-1}\).
\end{lemma}

\begin{proof}
First fix \(U\).  For arrows
\(r:B'_0\to B'_1\) and \(s:B_0\to B_1\) in \(\mathbb B(U)\), define
\[
  K_U(r^{\op},s):K_U(B'_1,B_0)\longrightarrow K_U(B'_0,B_1)
\]
as the unique map induced by the inner cowedge whose \((w,A)\)-component is
\begin{align*}
  H_W(A,A,w^*B'_1,w^*B_0)
  &\xrightarrow{H_W(1,1,(w^*r)^{\op},w^*s)}
  H_W(A,A,w^*B'_0,w^*B_1)
  \\
  &\xrightarrow{\kappa^{U,B'_0,B_1}_{w,A}}
  w^*K_U(B'_0,B_1).
\end{align*}
For \(z:Z\to W\), the base-coherence equation for this candidate is the
base-coherence equation for \(\kappa^{U,B'_0,B_1}\), pasted with the indexed
naturality square
\[
 \chi^H_z\,
 z^*H_W(1,1,(w^*r)^{\op},w^*s)
 =
 H_Z(1,1,((wz)^*r)^{\op},(wz)^*s)\,
 \chi^H_z,
\]
where the parameter \(\alpha\)'s are suppressed.  Thus indexed naturality
of \(H\), not fibrewise naturality alone, shows that this is an inner
relative cowedge.  Two successive such maps and the map associated to the
composite have identical
composites with every universal component \(\kappa_{w,A}\), so uniqueness in
the inner universal property makes them equal.  The same argument with
identity arrows proves the identity law.  Thus each \(K_U\) is a bifunctor
with the asserted variance.

For \(u:U'\to U\), set
\[
  \chi^K_{u,B',B}=(b_u^{B',B})^{-1}:
  u^*K_U(B',B)\xrightarrow{\ \cong\ }
  K_{U'}(u^*B',u^*B).
\]
For \(r:B'_0\to B'_1\) and \(s:B_0\to B_1\), naturality of the base map is
the equality
\begin{equation}
\begin{split}
 u^*K_U(r^{\op},s)\,b_u^{B'_1,B_0}
 ={}&
 b_u^{B'_0,B_1}\,
 K_{U'}\bigl((u^*r)^{\op},u^*s\bigr).
\end{split}
\label{eq:K-indexed-naturality}
\end{equation}
Precompose both sides with every \(\kappa^{U'}_{w,A}\) and apply
\cref{eq:kappa-base-change}.  After cancelling the same \(\alpha\)'s, the
two resulting maps differ only by whether the indexed naturality square of
\(H\) is traversed before or after
\(H(1,1,(w^*u^*r)^{\op},w^*u^*s)\); hence that square proves
\cref{eq:K-indexed-naturality}.  Universality of the inner coend then gives
the displayed equality.

For composable \(U''\xrightarrow{v}U'\xrightarrow{u}U\), put
\[
\begin{split}
 a'&=\alpha^{\mathbb B}_{v,u,B'}:
      v^*u^*B'\longrightarrow(uv)^*B',\\
 a&=\alpha^{\mathbb B}_{v,u,B}:
      v^*u^*B\longrightarrow(uv)^*B,
\end{split}
\]
and define the parameter transport
\begin{equation}
 \Theta_{v,u}^{B',B}
 =
 K_{U''}\bigl((a')^{\op},a^{-1}\bigr):
 K_{U''}((uv)^*B',(uv)^*B)
 \longrightarrow
 K_{U''}(v^*u^*B',v^*u^*B).
 \label{eq:K-parameter-transport}
\end{equation}
The correctly typed composition equation for the base maps is
\begin{equation}
 b_{uv}^{B',B}
 =
 \alpha^{\V}_{v,u,K_U(B',B)}\,
 v^*b_u^{B',B}\,
 b_v^{u^*B',u^*B}\,
 \Theta_{v,u}^{B',B}.
 \label{eq:K-base-composition}
\end{equation}
Indeed, precompose both sides with every universal component for
\(K_{U''}((uv)^*B',(uv)^*B)\).  Applying
\cref{eq:kappa-base-change} once to the left and successively for \(v\) and
\(u\) to the right reduces both sides to the same component of
\(\kappa^{U,B',B}\).  The two parameter transports agree by the pentagon
for \(\mathbb B\), the two apex transports agree by the pentagon for
\(\V\), and the intervening maps agree by the composition axiom for
\(\chi^H\).  Universality of the inner coend proves
\cref{eq:K-base-composition}.

For completeness, if
\(\iota^{\mathbb B}_{U,B}:1_U^*B\to B\) and
\(\iota^\V_{U,X}:1_U^*X\to X\) are the unitors, the corresponding typed
unit equation is
\begin{equation}
 \iota^\V_{U,K_U(B',B)}\,
 b_{1_U}^{B',B}\,
 K_U\bigl((\iota^{\mathbb B}_{U,B'})^{\op},
          (\iota^{\mathbb B}_{U,B})^{-1}\bigr)
 =1_{K_U(B',B)}.
 \label{eq:K-base-unit}
\end{equation}
It follows by the same one-component calculation, using the triangle
axioms instead of the pentagons.  Inverting
\cref{eq:K-base-composition} gives the explicit indexed-functor
composition law
\begin{equation}
\begin{split}
 \chi^K_{v,u^*B',u^*B}\,v^*\chi^K_{u,B',B}
 ={}&
 \Theta_{v,u}^{B',B}\,
 \chi^K_{uv,B',B}\,
 \alpha^\V_{v,u,K_U(B',B)}.
\end{split}
\label{eq:K-comparison-composition}
\end{equation}
\Cref{eq:K-base-unit} gives its unit law, and
\cref{eq:K-indexed-naturality} gives its naturality.  Therefore \(K\) is an
indexed functor.
\end{proof}

\begin{proposition}[Currying simultaneous cowedges]
\label{prop:cowedge-currying}
Under the hypotheses of \cref{lem:inner-indexed-functor}, there is a natural
bijection
\begin{equation}
  \RBiCow_U(H;X)
  \cong
  \RCow_U(K;X)
  \label{eq:cowedge-currying}
\end{equation}
for every \(U\) and \(X\in\V(U)\).
\end{proposition}

\begin{proof}
We construct both maps and verify every axiom.

\emph{From a simultaneous cowedge to an outer cowedge.}
Let \(\omega\) be simultaneous.  Fix \(v:V\to U\) and
\(B\in\mathbb B(V)\).  For \(w:W\to V\) and
\(A\in\mathbb A(W)\), define
\begin{equation}
  H_W(A,A,w^*B,w^*B)
  \xrightarrow{\omega_{vw,A,w^*B}}(vw)^*X
  \xrightarrow{\alpha^{-1}_{w,v,X}}w^*v^*X.
  \label{eq:omega-inner-cowedge}
\end{equation}
The \(A\)-dinaturality of \(\omega\) proves fibrewise dinaturality of this
family.  For \(z:Z\to W\), its base-coherence square is the joint
base-coherence square of \(\omega\), pasted with the associativity pentagon
for \(\alpha\).  Hence \cref{eq:omega-inner-cowedge} is an inner relative
cowedge over \(V\) with apex \(v^*X\).  Universality of \(K_V(B,B)\)
supplies a unique map
\begin{equation}
  \bar\omega_{v,B}:K_V(B,B)\longrightarrow v^*X
  \label{eq:baromega}
\end{equation}
whose pullback along every \(w\), after \(\kappa^{V,B,B}_{w,A}\), is
\cref{eq:omega-inner-cowedge}.

We check \(B\)-dinaturality.  For \(f:B\to C\) in \(\mathbb B(V)\), the
two required maps from \(K_V(C,B)\) to \(v^*X\) are
\[
  \bar\omega_{v,B}K_V(f^{\op},1_B),
  \qquad
  \bar\omega_{v,C}K_V(1_C,f).
\]
Precompose their pullbacks with every universal component of
\(K_V(C,B)\).  By the definition of the two actions of \(K\), the resulting
maps are respectively
\begin{align*}
  H_W(A,A,w^*C,w^*B)
  &\xrightarrow{H(1,1,(w^*f)^{\op},1)}
  H_W(A,A,w^*B,w^*B)
  \longrightarrow w^*v^*X,\\
  H_W(A,A,w^*C,w^*B)
  &\xrightarrow{H(1,1,1,w^*f)}
  H_W(A,A,w^*C,w^*C)
  \longrightarrow w^*v^*X.
\end{align*}
They agree by the \(B\)-dinaturality of \(\omega\).  The representing
property of the inner coend then makes the two maps from \(K_V(C,B)\)
equal.

For base coherence, let \(w:W\to V\).  We must show
\begin{equation}
  \bar\omega_{vw,w^*B}\,\chi^K_{w,B,B}
  =\alpha_{w,v,X}\,w^*\bar\omega_{v,B}
  :w^*K_V(B,B)\longrightarrow(vw)^*X.
  \label{eq:outer-base-check}
\end{equation}
Precompose both sides with
\(b_w^{B,B}:K_W(w^*B,w^*B)\to w^*K_V(B,B)\), which is an isomorphism.
Then precompose with every
\(\kappa^{W,w^*B,w^*B}_{z,A}\).  On the left,
\(\chi^K_{w,B,B}b_w^{B,B}=1\), so the defining property of
\(\bar\omega_{vw,w^*B}\) gives
\[
 H_Z(A,A,z^*w^*B,z^*w^*B)
 \xrightarrow{\omega_{(vw)z,A,z^*w^*B}}((vw)z)^*X
 \xrightarrow{\alpha^{-1}_{z,vw,X}}z^*(vw)^*X,
\]
with the parameter constraints understood.  On the right,
\cref{eq:kappa-base-change} converts
\(\kappa^W_{z,A}\) followed by \(z^*b_w\) into the
\((wz,A)\)-component of \(\kappa^V\).  The defining property of
\(\bar\omega_{v,B}\) then gives the same displayed map; the equality of the
two structural composites is the pentagon for \(\alpha\).
Universality of \(K_W(w^*B,w^*B)\), followed by invertibility of \(b_w\),
proves \cref{eq:outer-base-check}.  Hence
\(\bar\omega\) is an outer relative cowedge for \(K\).

\emph{From an outer cowedge to a simultaneous cowedge.}
Conversely, let
\(\tau\in\RCow_U(K;X)\).  Define
\begin{equation}
  \omega^\tau_{v,A,B}:
  H_V(A,A,B,B)
  \xrightarrow{\kappa^{V,B,B}_{1_V,A}}K_V(B,B)
  \xrightarrow{\tau_{v,B}}v^*X,
  \label{eq:uncurried-cowedge}
\end{equation}
with the unit constraints suppressed.  The inner cowedge equation for
\(\kappa\) gives \(A\)-dinaturality.  The definition of the two actions of
\(K\), followed by the \(B\)-dinaturality of \(\tau\), gives
\(B\)-dinaturality.  For \(w:W\to V\), the required compatibility is not
implicit: paste the inner base-coherence equation for the universal cowedge
\(\kappa\) with \cref{eq:kappa-base-change} for \(u=w\) and the identity
arrow of \(W\).  Postcomposing that equality with the outer
base-coherence equation for \(\tau\) gives
\[
  \omega^\tau_{vw,w^*A,w^*B}\chi^H_{w,A,A,B,B}
  =\alpha_{w,v,X}w^*\omega^\tau_{v,A,B}.
\]
Thus \(\omega^\tau\) is simultaneous.

\emph{The constructions are inverse.}
Starting with \(\omega\), the uncurried component at \((v,A,B)\) evaluates
the map \(\bar\omega_{v,B}\) on the identity-base universal component.
Its defining property \cref{eq:omega-inner-cowedge} therefore returns
\(\omega_{v,A,B}\), up to the unit constraint.  Starting with \(\tau\), the
inner cowedge used to construct \(\bar\omega_{v,B}\) from
\(\omega^\tau\) is the restriction of \(\tau_{v,B}\) along all
\(w:W\to V\).  More explicitly, \cref{eq:kappa-base-change}, together
with the base-coherence of the inner universal cowedge and the unit
constraints, gives
\begin{equation}
  \kappa^{W,w^*B,w^*B}_{1_W,A}
  =
  \chi^K_{w,B,B}\,
  \kappa^{V,B,B}_{w,A}.
  \label{eq:kappa-identity-restriction}
\end{equation}
The outer base-coherence axiom for \(\tau\) then yields
\begin{equation}
 \alpha^{-1}_{w,v,X}\,
 \tau_{vw,w^*B}\,
 \kappa^{W,w^*B,w^*B}_{1_W,A}
 =
 w^*\tau_{v,B}\,
 \kappa^{V,B,B}_{w,A}.
 \label{eq:tau-restricted-inner}
\end{equation}
Thus this inner cowedge is exactly the restriction of \(\tau_{v,B}\).
The uniqueness clause for the inner coend therefore gives
\(\bar\omega_{v,B}=\tau_{v,B}\).  Both assignments commute with
postcomposition in \(X\), proving naturality.
\end{proof}

\begin{theorem}[Relative Fubini]
\label{thm:relative-fubini}
Assume that the \(\mathbb A\)-inner coends
\(K_U(B',B)\) of \cref{eq:inner-K} exist and are stable under base change.
Then the simultaneous and iterated representing objects below have
equivalent existence: either both exist or neither does.  When they exist,
there is a canonical isomorphism
\begin{equation}
  \mathop{\int^{(A,B)\in\mathbb A\times\mathbb B}_{\Rel,U}}
     H(A,A,B,B)
  \ \cong\ 
  \mathop{\int^{B\in\mathbb B}_{\Rel,U}}K(B,B).
  \label{eq:relative-fubini}
\end{equation}
Here the right-hand side is the relative coend of the indexed bifunctor
\[
  K_V(B',B)
  =
  \mathop{\int^{A\in\mathbb A}_{\Rel,V}}H(A,A,B',B),
\]
not a list of unrelated fibrewise coends.

If, symmetrically, the \(\mathbb B\)-inner coends exist stably, then whenever
one of the simultaneous or iterated coends exists, all corresponding
two-variable admissible orders exist and are canonically isomorphic.
\end{theorem}

\begin{proof}
By \cref{lem:product-dinaturality,prop:cowedge-currying}, for every
\(X\in\V(U)\) there are natural bijections
\[
  \RCow_U(\widehat H;X)
  =\RBiCow_U(H;X)
  \cong\RCow_U(K;X).
\]
The simultaneous coend is, by \cref{prop:representing}, a representing
object for the first functor of \(X\); the iterated coend is a representing
object for the last.  The natural bijection therefore transports a universal
element in either direction.  Hence one object exists if and only if the
other does, and uniqueness of representing objects gives the canonical
isomorphism \cref{eq:relative-fubini}.

Repeating the proof with \(\mathbb A\) and \(\mathbb B\) interchanged gives
the other iterated order.  Every comparison is the unique isomorphism that
transports the same simultaneous universal cowedge.  No higher-variable
coherence is claimed here: it requires a separately defined multivariable
relative cowedge and recursive stability hypotheses for each intermediate
coend.
\end{proof}

\begin{remark}[Why stability is not cosmetic]
\label{rem:fubini-stability-typing}
If the inner objects \(K_U(B',B)\) merely exist separately, the canonical
maps in \cref{eq:K-base-map} need not be invertible.  Then they do not provide
the indexed comparison in the direction required for an indexed functor
\(K\), so the outer relative coend in \cref{eq:relative-fubini} is not typed
by the present definition.  The theorem therefore does not silently promote
fibrewise existence to indexed Fubini.  Under the hypotheses of
\cref{thm:stable-pullback}, the needed stability follows automatically from
the coequalizer presentation.
\end{remark}

\section{Further directions and present limits}
\label{sec:next}

Theorems \ref{thm:fixed-base-coequalizer}, \ref{thm:stable-pullback}, and
\ref{thm:relative-fubini} establish fixed-base existence, one checkable
stable range, and Fubini.  They do not yet identify the relative coend with
Shulman's canceling tensor, prove the relative co-Yoneda formula announced in
Drieux's slides, construct a profunctor equipment, or remove dependence on a
chosen site presentation.  The remaining problems are the following.

\begin{problem}[Comparison with indexed canceling products]
\label{prob:shulman-comparison}
When Shulman's hypotheses hold, determine whether the localized relative
coend is represented by his canceling tensor product.  State separately the
comparison on a fixed site and the additional descent needed for a stack of
coefficients.
\end{problem}

\begin{problem}[Relative co-Yoneda]
\label{prob:coyoneda}
For the indexed hom object suggested by \cref{eq:drieux-hom}, identify the
coefficient system and tensor needed to type the expression
\[
  \int^{A}_{\Rel} \D(A,X)\otimes F(A),
\]
then prove a stable equivalence with \(F(X)\).  The proof must distinguish
presheaf-level computation, sheafification, and descent.
\end{problem}

\begin{problem}[Relative profunctors]
\label{prob:relative-prof}
Use the existence and Fubini theorems to define composition of
localization-sensitive indexed profunctors.  Prove associativity and units,
and compare the result with Shulman's indexed profunctor equipment and Lee's
2-categorical construction \cite{Lee2023}.
\end{problem}

\begin{problem}[Site-presentation invariance]
\label{prob:site-invariance}
Let two small or small-generated sites present equivalent base topoi.  Find
the additional comparison data on \(\D,\V,H\) for which their relative-coend
objects correspond.  A satisfactory theorem should be invariant under the
relevant site Morita equivalence, not merely under a fibrewise equivalence
chosen on one presentation.
\end{problem}

These boundaries separate the present fixed-site calculus from the stronger
relative-topos and Morita-invariant theory that remains to be developed.
